% Generated draft: requires mathematical and source review.
\documentclass[11pt]{article}
\usepackage[margin=25mm]{geometry}
\usepackage{fontspec}
\setmainfont{Latin Modern Roman}
\usepackage{amsmath,amssymb,mathtools,unicode-math}
\setmathfont{Latin Modern Math}
\usepackage{enumitem,needspace,xurl,hyperref}
\setlength{\parindent}{0pt}
\setlength{\parskip}{5pt}
\setlength{\emergencystretch}{3em}
\newcommand{\N}{\mathbb{N}}
\newcommand{\Z}{\mathbb{Z}}
\newcommand{\Q}{\mathbb{Q}}
\newcommand{\R}{\mathbb{R}}
\newcommand{\C}{\mathbb{C}}
\newcommand{\F}{\mathbb{F}}
\newcommand{\A}{\mathbb{A}}
\newcommand{\PP}{\mathbb{P}}
\newcommand{\E}{\mathbb{E}}
\newcommand{\eps}{\varepsilon}
\newcommand{\dd}{\,\mathrm{d}}
\begin{document}
\section*{UnsolvedMath ID: 11000203}
\textbf{Source problem code:} catalogue entry
\subsection{Definitions and mathematical statement}
Let \(S\) and \(M_{g,m}\) be the source's surface and mapping class group, and distinguish the missing preceding predicate \(\mathrm{CC}_{\mathrm{src}}\) from Farb--Mosher convex cocompactness \(\mathrm{CC}_{\mathrm{FM}}\). The compound problem asks:
\begin{enumerate}
\item Do these two predicates agree for subgroups of \(M_{g,m}\)? The preceding source definition is absent and is not replaced by an assumed curve-graph or Teichm\"uller criterion.
\item For each such convex cocompact subgroup \(\Gamma\), is its natural surface-group extension
\[
 1\longrightarrow\pi_1(S)\longrightarrow E_\Gamma
 \longrightarrow\Gamma\longrightarrow1
\]
word hyperbolic, meaning a finitely generated group with a thin-triangle Cayley graph? The intended extension is induced by the mapping-class outer action, but its construction and closed versus punctured surface hypotheses must be supplied. Peripheral subgroups make that distinction material; no relative-hyperbolicity conclusion is substituted.
\item Is there a convex cocompact subgroup \(\Gamma\leq M_{g,m}\) isomorphic to \(\pi_1(\Sigma_h)\) for a closed surface with \(h\geq2\)?
\end{enumerate}

The accompanying discussion uses ``Veech surfaces'' for projections of complex Teichm\"uller geodesic disks whose stabilizers are lattices in \(\operatorname{PSL}(2,\R)\). In this usage they are curves/orbifold surfaces in moduli space, not the original flat surface. The disk is described as maximal embedded in a bounded-domain model in \(\C^{3g-3+m}\). The source reports these quotients as finite type with isolated singularities and never closed [V], so their corresponding groups are virtually free, with distinguished conjugacy classes from cusps. It cites [McM1, McM2] for properties, [L] for elementary constructions and Leininger--Reid [LR] for explicit subgroup constructions. These contextual results are not asserted to produce the closed-surface subgroup in item (3).
\subsection{Short English statement}
Compare the two convex-cocompactness definitions, ask whether the associated surface-group extension is hyperbolic, and seek a convex cocompact closed-surface subgroup, retaining the Veech-curve context.
\subsection{Sources}
\begin{itemize}
\item \emph{Problems on Mapping Class Groups and Related Topics}, edited by Farb, Problem 4, supplied PDF page 247. \url{https://www.math.uchicago.edu/~farb/papers/mcgbook.pdf}. Recorded access 2026-07-29; extraction pdf-font-structure.
\item Farb--Mosher [FMo]; [V] for nonclosedness of Veech quotients; [McM1], [McM2], [L] and Leininger--Reid [LR] as described in the source. Complete bibliography entries are absent. The summary mentions later Hamenst\"adt and other equivalence work without a precise theorem.
\item Local UnsolvedMath registry, \texttt{data/unsolvedmath/problems.json}, permanent ID \(11000203\), exact source code \texttt{catalogue entry}; supplied OPEN-TRIAGE summary tentatively reports partial equivalence progress.
\end{itemize}
\end{document}
